\section{Inverse comorbidity：跨领域连接如何交给人类}

最后一个主讲案例不从已知机制开始，而从一个反常统计关系开始：某些癌症与神经退行性疾病之间可能存在 inverse comorbidity，即一个疾病出现时，另一个疾病的发生率低于预期。系统尝试寻找保护癌细胞的机制是否会伤害神经元，或反过来。

\subsection{什么是 inverse comorbidity}

普通 comorbidity 指两种疾病共同出现超过独立预期；inverse comorbidity 则低于预期。若疾病 $A$ 与 $B$ 在独立假设下的共同概率为 $p(A)p(B)$，可定义一个简单比率：
$$
\operatorname{ICR}(A,B)=\frac{p(A\cap B)}{p(A)p(B)}.
$$
其中：
\begin{itemize}
  \item $p(A\cap B)$ 是两种疾病共同出现的观察概率；
  \item $p(A)p(B)$ 是独立时的期望共同概率；
  \item $\operatorname{ICR}<1$ 提示 inverse association；
  \item 年龄、死亡竞争、诊断偏差和治疗差异都可能制造假关联。
\end{itemize}

\begin{knowledgebox}{背景概念：inverse comorbidity}
逆共病是提出机制问题的入口，不是因果结论。若 SCLC 患者较少出现某类神经退行性疾病，可能是癌细胞激活了保护性程序，也可能只是生存期、年龄结构或医疗记录造成。好的 co-scientist 应同时生成生物机制与流行病学反解释。
\end{knowledgebox}

\lecturefigure{slide-020.jpg}{SCLC 与 neurodegeneration 的 inverse-comorbidity 案例}{Stanford CS25 V6 Lecture 07 官方录像 00:57:24；课堂未发表材料}

\begin{warningbox}{案例状态}
本节展示的是系统生成的研究报告和一次专家联系，不是已完成的实验论文。它的教学价值在于跨领域 hypothesis generation 与 human handoff，而不在于证明 SCLC 机制可以治疗神经退行性疾病。
\end{warningbox}

\subsection{把反常关系变成可研究问题}

问题被具体化为：small cell lung cancer（SCLC）细胞如何在高突变、基因组混乱和死亡压力下存活？这些生存机制是否与神经元在 Alzheimer’s 等疾病中失败的路径形成镜像？系统需要连接癌症 genomics、RNA editing、cell death、micro-exon 和 neuronal maintenance，多数线索原本分散在不同文献社区。

\lecturefigure{slide-021.jpg}{SCLC 与 neurodegeneration 的研究问题}{Stanford CS25 V6 Lecture 07 官方录像 00:58:02}

\begin{knowledgebox}{读图：先区分统计观察与生物学目标}
左侧 bullet 说明癌细胞受到基因组混乱却持续生长，神经元则在疾病中逐步死亡；inverse comorbidity 只是提示两种命运可能共享反向机制。右侧细胞图是研究对象视觉，不是因果证据。真正的研究目标是找到可测量的保护程序，并设计能区分多种解释的实验。
\end{knowledgebox}

\subsection{四类候选：从 genomic chaos 到 neural hijack}

系统给出多个方向，而不是只押注一个漂亮故事。Genomic chaos 假设聚焦 ADAR1 等 RNA editing 防御；DHX9 与 SRRM4 方向连接 helicase 和 neuronal micro-exon；Evading death 研究癌细胞怎样抑制死亡；Neural hijack 则追问肿瘤是否重用神经发育程序。Portfolio 的价值在于让实验者比较互斥预测。

\lecturefigure{slide-022.jpg}{连接 SCLC 与 neurogenesis 的候选假设组合}{Stanford CS25 V6 Lecture 07 官方录像 00:58:18}

\begin{knowledgebox}{读图：不要把四张卡片读成四个事实}
每张卡片是一个搜索方向：机制、关键分子、可联系专家和潜在实验都不同。系统应保留它们之间的竞争关系，例如 ADAR1/DHX9 的“暗基因组屏蔽”解释与神经 micro-exon 解释可能给出不同干预预测。好的 meta-review 不应过早压缩成唯一故事。
\end{knowledgebox}

\begin{knowledgebox}{术语消化：候选机制}
\begin{tabularx}{\textwidth}{L{0.22\textwidth}X X}
\toprule
方向 & 解决的问题 & 可验证线索 \\
\midrule
Genomic chaos & 癌细胞怎样承受异常 RNA 与突变压力 & 干预 ADAR1、观察 Z-RNA 与死亡信号 \\
DHX9/SRRM4 & RNA helicase 与 neuronal micro-exon 是否形成保护差异 & 比较表达、剪接与存活表型 \\
Evading death & 癌细胞抑制哪类 programmed cell death & 激活或阻断候选通路 \\
Neural hijack & 肿瘤是否重用神经发育或维持程序 & 追踪神经标记、转录网络与功能 \\
\bottomrule
\end{tabularx}
\end{knowledgebox}

\subsection{Cold email 能证明什么}

研究者把系统提出的候选交给领域专家 Dr. Bellegia，得到“extreme outlier”式的积极反馈：相关基因在 SCLC 数据中异常突出。这个反馈提高了候选的外部 plausibility，也证明系统能找到合适联系人，但专家认可仍可能受选择、叙事和已知证据影响。

\lecturefigure{slide-023.jpg}{领域专家对候选基因异常性的反馈}{Stanford CS25 V6 Lecture 07 官方录像 00:58:50}

\begin{knowledgebox}{读图：排名高不等于机制已验证}
条形图把 SCLC 中的候选基因表达放在其他癌症背景中，显示其为显著 outlier。它支持“这个方向值得进一步研究”，却没有证明该基因导致 inverse comorbidity。下一步应有预注册的细胞或动物干预、独立数据集和针对替代解释的 negative control。
\end{knowledgebox}

\begin{warningbox}{Expert email 的证据位置}
专家反馈高于完全内部自评，因为它引入外部领域判断；但它低于盲评、正式 peer review 和实验复现。公开演讲最容易只展示回复积极的邮件，因此应记录联系了多少专家、无回复与否定意见，以及候选在联系前的冻结排名。
\end{warningbox}

\teachervoice{00:57:24--01:00:02，老师把这一段称为 decoding inverse comorbidity，并展示系统怎样给出研究联系人。课堂强调的是 unexpected connection 和协作启动，不是已完成的生物学突破。}

\subsection{百页报告怎样避免淹没科学家}

系统最终报告包含候选机制、图示、证据、反例和研究联系人。长报告的优点是保留低排名线索，缺点是阅读成本与“完整性幻觉”：格式越像正式论文，用户越容易忽略其本质仍是候选集合。报告必须让 provenance、uncertainty、rank history 和建议实验可单独审计。

\lecturefigure{slide-024.jpg}{AI co-scientist 生成的详细候选机制报告}{Stanford CS25 V6 Lecture 07 官方录像 00:59:22}

\begin{knowledgebox}{读图：图文完整度不是证据完整度}
页面给出 ADAR1/Z-RNA 假设、机制示意和实验建议，看起来像成熟论文。读者应逐项检查：引用是否支持具体边；图是实验数据、公开图还是生成示意；哪些句子是已知事实，哪些是推断；候选为何排第一；哪些结果会推翻它。报告的最佳用途是压缩研究启动成本，而不是绕过审查。
\end{knowledgebox}

\lecturefigure{slide-025.jpg}{报告中的 unexpected connections 与建议研究联系人}{Stanford CS25 V6 Lecture 07 官方录像 00:59:38}

\begin{knowledgebox}{读图：从机制组合过渡到 human handoff}
上半部分把 ADAR1、DHX9、p53、SRRM4 等线索组合成 sequential shield collapse；下半部分列出可联系的领域专家及其知识位置。这里体现 Agent 的两个价值：发现文献社区之间的连接，并把下一步验证路由给具体人。风险是联系人和机制同样可能由不完整检索产生，需要人工核验身份、贡献与利益冲突。
\end{knowledgebox}

\begin{lstlisting}[language=Python,caption={长报告的可审计交付结构}]
report = {
    "hypotheses": ranked_hypotheses,
    "edge_level_sources": provenance_graph,
    "counterevidence": reflection.failures,
    "proposed_experiments": falsification_tests,
    "uncertainty": calibrated_confidence,
    "contacts": verified_domain_experts,
    "run_metadata": model_tools_budget_and_timestamps,
}
\end{lstlisting}

\subsection{本章小结}

Inverse comorbidity 案例展示了 AI co-scientist 的跨领域搜索价值：从统计反常出发，形成多个竞争机制，找到外部专家，再输出带研究联系人的详细报告。它仍处于假设和人类交接阶段。正确的下一步不是宣传单一故事，而是冻结候选、设计能区分机制的实验，并记录负面专家反馈和失败路线。

